\section{Related Work}
\label{sec:related}

Definitions of verifiability separate goals from judges. Non-interactive
verifiable computation addresses outsourcing \cite{gennaro2010},
proof-carrying data distributes a compliance predicate \cite{chiesa2010pcd},
and Cortier et al. parameterize voting verifiability by a goal, judge, tolerance,
and benign-run condition \cite{cortier2016}. We instantiate that scaffold for
one fusion run.

Verifiable aggregation protocols instead bundle assurances. SIA authenticates
reports, commits to values, and checks
operator-specific approximations with sampling and interactive proofs
\cite{przydatek2003sia}. Hierarchical secure aggregation prevents intermediate
manipulation beyond the values injected directly by compromised nodes
\cite{chan2006hierarchical}. Recent schemes add public verification for private
summation \cite{palazzo2024mpvas} or designated verification of private inner
products \cite{zhou2024ppvda}. Their theorems are packages for particular
operators and adversaries, not cumulative names for the run-level conditions
inside those packages.

Mechanism surveys classify protocol families rather than claims.
Bontekoe et al. \cite{bontekoe2025} organize verifiable privacy-preserving
computation by computation and verification technique; our map instead puts
the atomic audit condition first.

Provenance and runtime assurance protect adjacent objects: SLSA orders build
evidence \cite{slsa}, RATS appraises an attesting environment \cite{rfc9334},
and APEX proves execution under software compromise \cite{apex2020}. None is
itself a cumulative V-level.

Robust fusion addresses what cryptography cannot validate. MATE infers which
sensors to believe from cross-sensor consistency \cite{mate2025}; V-goals
instead check bound records and computation. Privacy is likewise orthogonal
unless its mechanism supplies an audit condition.
